% GENERATED FILE. Edit canonical lessons, then run npm run generate:books.
% canonical-source-hash: fnv1a64:9263ddbae9b4e0f5
% canonical-lessons: KA-C123-tarakari, KA-C123-balehannu, KA-C123-mavinahannu, KA-C123-tengu, KA-C123-irulli

\chapter{Vegetables, Banana, Mango, Coconut, Onion}
\label{ch:ka-vegetables-banana-mango-coconut-onion}
\hlchaptermodality{123}

\begin{chapteropening}
\textbf{By the end of this chapter:} \emph{I can say vegetables, a banana, a mango, a coconut and an onion.}

Complete the last lesson of chapter 123: I can say vegetables, a banana, a mango, a coconut and an onion.
\end{chapteropening}

\section[\texorpdfstring{\kn{ತರಕಾರಿ} (tarakāri)}{ತರಕಾರಿ (tarakāri)}]{\kn{ತರಕಾರಿ} (tarakāri) — vegetables}

\label{lesson:KA-C123-tarakari}



\begin{quote}
Before the new one: say the Kannada for a fish, then the Kannada for meat.
\end{quote}

\subsection*{You'll want to know: \kn{ತರಕಾರಿ}}
\textbf{\kn{ತರಕಾರಿ}} — \emph{tarakāri} — \textquotedblleft{}vegetables\textquotedblright{}.

\begin{grammarlens}[title={What you've built}]
One more word for this chapter.
\end{grammarlens}

\subsection*{Your turn}
\begin{itemize}
\raggedright
  \item \emph{Say it:} \emph{tarakāri}
  \item \emph{Say it:} \emph{tarakāri}, once more
  \item \emph{Say it:} say \emph{tarakāri}
  \item \emph{Recall:} say the Kannada for a fish, then the Kannada for meat, then say \emph{tarakāri} again
\end{itemize}

\begin{tcolorbox}[breakable,colback=teal!4,colframe=teal!35!black,title={Before you move on}]
What does \kn{ತರಕಾರಿ} mean? (\textquotedblleft{}Vegetables\textquotedblright{}.) Say it once more.
\end{tcolorbox}
\section[\texorpdfstring{\kn{ಬಾಳೆಹಣ್ಣು} (bāḷehaṇṇu)}{ಬಾಳೆಹಣ್ಣು (bāḷehaṇṇu)}]{\kn{ಬಾಳೆಹಣ್ಣು} (bāḷehaṇṇu) — a banana}

\label{lesson:KA-C123-balehannu}



\begin{quote}
Before the new one: say the Kannada for meat, then the Kannada for vegetables.
\end{quote}

\subsection*{You'll want to know: \kn{ಬಾಳೆಹಣ್ಣು}}
\textbf{\kn{ಬಾಳೆಹಣ್ಣು}} — \emph{bāḷehaṇṇu} — \textquotedblleft{}a banana\textquotedblright{}.

\begin{grammarlens}[title={What you've built}]
One more word for this chapter.
\end{grammarlens}

\subsection*{Your turn}
\begin{itemize}
\raggedright
  \item \emph{Say it:} \emph{bāḷehaṇṇu}
  \item \emph{Say it:} \emph{bāḷehaṇṇu}, once more
  \item \emph{Say it:} say \emph{bāḷehaṇṇu}
  \item \emph{Recall:} say the Kannada for meat, then the Kannada for vegetables, then say \emph{bāḷehaṇṇu} again
\end{itemize}

\begin{tcolorbox}[breakable,colback=teal!4,colframe=teal!35!black,title={Before you move on}]
What does \kn{ಬಾಳೆಹಣ್ಣು} mean? (\textquotedblleft{}A banana\textquotedblright{}.) Say it once more.
\end{tcolorbox}
\section[\texorpdfstring{\kn{ಮಾವಿನಹಣ್ಣು} (māvinahaṇṇu)}{ಮಾವಿನಹಣ್ಣು (māvinahaṇṇu)}]{\kn{ಮಾವಿನಹಣ್ಣು} (māvinahaṇṇu) — a mango}

\label{lesson:KA-C123-mavinahannu}



\begin{quote}
Before the new one: say the Kannada for vegetables, then the Kannada for a banana.
\end{quote}

\subsection*{You'll want to know: \kn{ಮಾವಿನಹಣ್ಣು}}
\textbf{\kn{ಮಾವಿನಹಣ್ಣು}} — \emph{māvinahaṇṇu} — \textquotedblleft{}a mango\textquotedblright{}.

\begin{grammarlens}[title={What you've built}]
One more word for this chapter.
\end{grammarlens}

\subsection*{Your turn}
\begin{itemize}
\raggedright
  \item \emph{Say it:} \emph{māvinahaṇṇu}
  \item \emph{Say it:} \emph{māvinahaṇṇu}, once more
  \item \emph{Say it:} say \emph{māvinahaṇṇu}
  \item \emph{Recall:} say the Kannada for vegetables, then the Kannada for a banana, then say \emph{māvinahaṇṇu} again
\end{itemize}

\begin{tcolorbox}[breakable,colback=teal!4,colframe=teal!35!black,title={Before you move on}]
What does \kn{ಮಾವಿನಹಣ್ಣು} mean? (\textquotedblleft{}A mango\textquotedblright{}.) Say it once more.
\end{tcolorbox}
\section[\texorpdfstring{\kn{ತೆಂಗು} (teṅgu)}{ತೆಂಗು (teṅgu)}]{\kn{ತೆಂಗು} (teṅgu) — a coconut}

\label{lesson:KA-C123-tengu}



\begin{quote}
Before the new one: say the Kannada for a banana, then the Kannada for a mango.
\end{quote}

\subsection*{You'll want to know: \kn{ತೆಂಗು}}
\textbf{\kn{ತೆಂಗು}} — \emph{teṅgu} — \textquotedblleft{}a coconut\textquotedblright{}.

\begin{grammarlens}[title={What you've built}]
One more word for this chapter.
\end{grammarlens}

\subsection*{Your turn}
\begin{itemize}
\raggedright
  \item \emph{Say it:} \emph{teṅgu}
  \item \emph{Say it:} \emph{teṅgu}, once more
  \item \emph{Say it:} say \emph{teṅgu}
  \item \emph{Recall:} say the Kannada for a banana, then the Kannada for a mango, then say \emph{teṅgu} again
\end{itemize}

\begin{tcolorbox}[breakable,colback=teal!4,colframe=teal!35!black,title={Before you move on}]
What does \kn{ತೆಂಗು} mean? (\textquotedblleft{}A coconut\textquotedblright{}.) Say it once more.
\end{tcolorbox}
\section[\texorpdfstring{\kn{ಈರುಳ್ಳಿ} (īruḷḷi)}{ಈರುಳ್ಳಿ (īruḷḷi)}]{\kn{ಈರುಳ್ಳಿ} (īruḷḷi) — an onion}

\label{lesson:KA-C123-irulli}



\begin{quote}
Before the new one: say the Kannada for a mango, then the Kannada for a coconut.
\end{quote}

\subsection*{You'll want to know: \kn{ಈರುಳ್ಳಿ}}
\textbf{\kn{ಈರುಳ್ಳಿ}} — \emph{īruḷḷi} — \textquotedblleft{}an onion\textquotedblright{}.

\begin{grammarlens}[title={What you've built}]
One more word for this chapter.
\end{grammarlens}

\subsection*{Your turn}
\begin{itemize}
\raggedright
  \item \emph{Say it:} \emph{īruḷḷi}
  \item \emph{Say it:} \emph{īruḷḷi}, once more
  \item \emph{Say it:} say \emph{īruḷḷi}
  \item \emph{Recall:} say the Kannada for a mango, then the Kannada for a coconut, then say \emph{īruḷḷi} again
\end{itemize}

\begin{tcolorbox}[breakable,colback=teal!4,colframe=teal!35!black,title={Before you move on}]
What does \kn{ಈರುಳ್ಳಿ} mean? (\textquotedblleft{}An onion\textquotedblright{}.) Say it once more.
\end{tcolorbox}
